% Part A (non-tech core) of the AI-for-X seminars: what AI and generative AI are. 4 frames.
% Adapted from ai_educators_content.tex (Big Picture, Language Model, Large). Audience-neutral on purpose.
% Part A is shared: edit here, then recompile every deck that includes it. Trim by commenting frames, not deleting.
% Size target: same as ai_core_tech_whatisai (4 frames).

%%%%%%%%%%%%%%%%%%%%%%%%%%%%%%%%%%%%%%%%%%%%%%%%%%%%%%%%%%%
\begin{frame}[fragile]\frametitle{The Big Picture: AI, ML, and Deep Learning}
\begin{center}
\includegraphics[width=0.55\linewidth,keepaspectratio]{ai1}
\end{center}
\begin{itemize}
\item \textbf{AI}: any machine doing something that looks intelligent
\item \textbf{Machine Learning (ML)}: a machine that learns the pattern from data, instead of being told the rule directly
\item \textbf{Deep Learning (DL)}: ML using many-layered neural networks, powerful enough to learn something as complex as language
\end{itemize}
{\tiny (Ref: https://blogs.nvidia.com/blog/2016/07/29/whats-difference-artificial-intelligence-machine-learning-deep-learning-ai/)}
\end{frame}

%%%%%%%%%%%%%%%%%%%%%%%%%%%%%%%%%%%%%%%%%%%%%%%%%%%%%%%%%%%
\begin{frame}[fragile]\frametitle{What Is a Language Model?}
\begin{itemize}
\item While typing an SMS, have you noticed it suggests the next word?
\item Those suggestions are not random: the phone has learned what word usually comes next from millions of past sentences.
\item Do this once, word by word, and you get a suggestion. Do it again and again, and you get a full sentence.
\item That is a \textbf{Language Model}, predicting the next word. Chained together, it writes whole paragraphs: a \textbf{Generative Model}.
\end{itemize}
\begin{center}
\includegraphics[width=0.4\linewidth,keepaspectratio]{chatgpt34}
\end{center}
\end{frame}

%%%%%%%%%%%%%%%%%%%%%%%%%%%%%%%%%%%%%%%%%%%%%%%%%%%%%%%%%%%
\begin{frame}[fragile]\frametitle{Why Are They Called ``Large''?}
\begin{itemize}
\item ChatGPT works on exactly this idea, just at an enormous scale: trained on a large fraction of everything ever written on the internet.
\item That is why the phrase is \textbf{Large Language Model (LLM)}: large training data, large number of internal parameters.
\item More data and more parameters mean better predictions, but the underlying idea never changes: predict the next word, one at a time.
\end{itemize}
\begin{center}
\includegraphics[width=0.5\linewidth,keepaspectratio]{chatgpt31}
\end{center}
\end{frame}

%%%%%%%%%%%%%%%%%%%%%%%%%%%%%%%%%%%%%%%%%%%%%%%%%%%%%%%%%%%
\begin{frame}[fragile]\frametitle{Generative AI: What It Can Make}
\begin{itemize}
\item \textbf{Generative AI} makes new content from a request: text, images, code, audio, even short video.
\item The chat assistants you may have met (ChatGPT, Claude, Gemini and others) are built on language models like the one on the last slide.
\item Same idea, different material: the model has seen very many examples of the thing, and produces a new one that looks like them.
\item It is a very fast, very well-read assistant. It is not a person, and it does not understand your situation unless you tell it.
\end{itemize}
\end{frame}
